\subsection{The actual quadratic source correction in the last compact orbit}

\begin{statement}{D.66}{Proved}{Two literal peripheral corrections fill the remaining old compact comparison block}
Retain the actual sources, arithmetic element, coefficient modules and
marked quotient maps of D.65. In particular,
\[
 X=\widehat B/(R_{50}\cap L_2),\qquad
 Y=X/\mathcal K,\qquad X_{65}=X/\mathcal J_{65},\qquad
 S_{57}=\mathfrak R_h\cap M_*\cap L_{*,2}.
 \tag{D.66.1}
\]
The integral homology image of \(S_{57}\) is the actual lattice
\(\mathcal L=\ker(U^-\to U_*^-)\). The coefficient group remains
the actual \(G\) of order \(576\cdot3^{97}\), with its original
source generators \(s,r\) and cusp \(t=sr\) of order 36.

There is an actual continuous peripheral automorphism of
\[
 \begin{aligned}
 X_{66}&=X/\mathcal J_{66},\\
 \mathcal J_{66}
  &=\ker\bigl(\mathcal J_{63}
        \xrightarrow{(q_1,q_2,q_3,q_4)}
                         C_1\oplus C_2\oplus C_3\oplus C_4\bigr).
 \end{aligned}
 \tag{D.66.2}
\]
It fixes the original cusp, has multiplier one, is identity on \(G\),
and retains the complete literal D.57 small action and the prescribed
\(\mathcal W_h\) action. All the proved normalized parameter,
metabelian and CM comparisons through \(\mathcal W_h\) are retained.
Both arithmetic branches and every retained small framing are included.
The homology action is the replacement described below.

The exact new projection kernel and dimensions are
\[
 \begin{aligned}
 \ker(X_{66}\to X_{65})
     &=\mathcal J_{65}/\mathcal J_{66}\simeq\ker\nu_4,\\
 \dim\ker\nu_4&=1918,
       &\dim(\ker\nu_4)_G&=4,\\
 \dim(\mathcal J_{63}/\mathcal J_{66})&=9600,
       &\dim(\mathcal J_{63}/\mathcal J_{66})_G&=28,\\
 \dim\mathcal J_{66}&=\dim\mathcal J_{65}-1918,\\
 \dim(\mathcal J_{66})_{\text{specified small characters}}&=274560.
 \end{aligned}
 \tag{D.66.3}
\]
Here \(\nu_4\) is the actual full norm used in D.65. The
four-dimensional coinvariant space of its kernel maps with rank two
to the coinvariants of \(C_4\); both maps are retained.

We give the finite source detecting the new correction. On the old
orbit \(a=4\), let \(\mathsf M\) be the integral lattice induced
from the actual rank-16 \(\mathcal O=\mathbb Z_2[\omega]\)-cover
through the eight unordered character cosets. Its \(\mathbb Z_2\)-rank
is 256. The actual small compact lattice \(\mathsf M_*\) has rank 64.
Write \(\phi:\mathsf M\twoheadrightarrow\mathsf M_*\) for their
marked map. There are actual marked epimorphisms
\[
 \begin{gathered}
 X\longrightarrow \mathscr E
       \longrightarrow\mathscr E_*,\\
 \mathscr E=\mathscr N\rtimes G^\dagger,\quad
 |G^\dagger|=46656,\quad
 \mathscr N^{\rm ab}=\mathsf M/4,\quad
 \mathscr N'=C\simeq\mathbb F_2^8,\\
 \mathscr E_*=\mathscr N_*\rtimes G_*,\quad
 |G_*|=3888,\quad
 \mathscr N_*^{\rm ab}=\mathsf M_*/4,\quad
 \mathscr N_*'=C_*\simeq\mathbb F_2^4 .
 \end{gathered}
 \tag{D.66.4}
\]
Both two-groups have class two and exponent four, and orders
\(2^{520}\) and \(2^{132}\). The second map induces the original
base map and \(\phi\). Its central map has the retained columns
\[
 (\mathtt{2},\mathtt{1},\mathtt{4},\mathtt{2},
   \mathtt{8},\mathtt{8},\mathtt{1},\mathtt{4}).
 \tag{D.66.5}
\]
Typewriter masks denote hexadecimal binary-coordinate masks.
The image of \(S_{57}\) in \(\mathscr N\) is exactly
\(\ker(\mathscr N\to\mathscr N_*)\), with the integral
homology and central lifts specified in the proof.

In this actual marked quotient the source cusp has zero
\(\mathsf M/4\)-coordinate and central coordinate
\(c_t=\mathtt{26}\). Its base coordinate has order twelve.
For \(\ell(c)=c_4+c_5\), with zero-based indices,
\[
 \ell(c_t)=1,\qquad
 V_0=\langle\mathtt{09},\mathtt{42}\rangle
       \subset\ker(C\to C_*).
 \tag{D.66.6}
\]
The image of \(\mathcal J_{65}\) in \(C\) is exactly \(V_0\).
There are actual relative compact automorphisms whose central
actions are all four maps
\[
                    I+v\ell,\qquad v\in V_0.
 \tag{D.66.7}
\]
They fix \(\ker(C\to C_*)\) pointwise. Their two literal
conjugator differences have zero image in the full small
class-two quotient. Their cusp differences are respectively
\(\mathtt{00},\mathtt{09},\mathtt{42},\mathtt{4b}\).
The original source \(s\) has order two and \(r\) order three;
all four corrections have zero effect on the displayed image
of \(t^{18}\).

The earlier integral test has an exact limitation. The secondary
modulo-four norm on the two residual old compact coordinates is
zero, as is the induced even-homology divided-orientation image.
Both coordinates inject into the whole compact coinvariants.
The new lift uses (D.66.7) and its literal quadratic cusp equation.

The full remaining source kernel is
\[
 \begin{aligned}
 \ker(X\to X_{66})&=\mathcal J_{66},\qquad
 \ker(X_{66}\to Y)=\mathcal K/\mathcal J_{66},\\
 S_{66}&=\{u\in R_{50}:
         \operatorname{im}_{\mathcal T_2}(u)\in\mathcal J_{66}\},\\
 S_{66}/(R_{50}\cap L_2)&=\mathcal J_{66}.
 \end{aligned}
 \tag{D.66.8}
\]
Its specified remaining central block has positive dimension
\(274560\). The remaining arithmetic boundary and the full
profinite realization of D.13/D.19 are not determined here.
The manuscript passages are Part D(e), section 51.16, folio 734
verso, scan 174, and section 51.15, folios 725--727, scans 156--160.
\end{statement}
\begin{proof}
\emph{Integral marked projective lattices.}
Use the old character \((0,1,0,0)\), its unordered stabilizer
\(H\) of order 72, and the actual graph of D.60. The affine,
cusp and compact modules on this orbit are projective over
the actual integral coefficient group, by D.60.12 and its
integral lifting argument. The old compact cover is a
projective direct summand. The small compact module is
projective over its actual small coefficient group as well.

The program \path{d_old_compact_integral.py} retains the D.64
cycles modulo four and their original reduction. In the real
32-dimensional cover module it chooses four generators for
the free Sylow-eight restriction. The 32 translated columns
form a basis. The coefficient-of-identity rows give an
ordinary map whose Sylow trace is the identity. Averaging
over the odd scalar group and the remaining stabilizer
cosets lifts the retained projector. The polynomial
\(3e^2-2e^3\) makes it idempotent modulo four. The rank-16
cover, all 72 semilinear actions, and the actual rank-four
small map retain their D.64 reductions. The small map kills
the rank-32 complementary projective summand exactly.

These maps lift integrally. Idempotents lift through
successive powers of two; projectivity lifts the prescribed
inclusion and retraction, as in D.61. The actual graph and
cusp kernels are saturated. Applying the projective-domain
Hom functor lifts their specified comparison maps and
reductions. The odd quotient from \(G\) to \(G^\dagger\)
is the actual D.60.10 map, with kernel of order \(3^{93}\).
Its normalized transfer is the inverse lattice comparison,
with odd denominators. Thus \(\mathsf M\), its graph
projection and its small map are actual integral summands
and maps of the full arithmetic homology.

\emph{The complete compact central comparisons.}
On one unordered pair the odd-invariant exterior module is
the D.64 Hermitian module. In the real basis \(u_i,\omega u_i\),
use exactly the diagonal and off-diagonal exterior
identifications of D.64.12. Its complete stabilizer
coinvariant quotient has dimension eight. The actual small
Hermitian quotient has dimension four. Applying
\(B\mapsto\phi B\bar\phi^{\,t}\) to all 256 Hermitian
basis vectors gives (D.66.5).

The comparison-kernel coinvariants inject into the whole
coinvariant quotient, with columns
\[
       \mathtt{09},\quad\mathtt{30},\quad
       \mathtt{42},\quad\mathtt{84}.
 \tag{D.66.9}
\]
They are exactly the kernel of (D.66.5). The ordinary norm
has rank two on this space; its kernel there is \(V_0\).
This verifies the central comparison through the actual
exterior map, rather than a dimension count alone.

The full induced \(Q\)-module has dimension 1920.
The program \path{d_old_compact_global_kernel.py} constructs
its source \(s,r\) matrices through the same eight coset
representatives and complete stabilizer actions. It checks
both generator relations. The full norm has rank two.
An actual basis of its 1918-dimensional kernel gives 3836
augmentation columns of rank 1914. Its coinvariant
dimension is therefore four. Its further map into the
eight-dimensional whole compact quotient has rank two
and image exactly \(V_0\). The odd parameter subgroup
acts trivially on the averaged self module, with odd norm
factor. These are also the actual \(G\)-module maps.

The secondary integral test is kept separately. Lift the
240-dimensional comparison kernel into the actual integral
Hermitian kernel modulo four by subtracting a right-inverse
image for \(\phi\). On an ordinary-norm-kernel vector,
divide the full signed integral norm by two, then reduce
modulo the ordinary norm image. A changed lift adds an
ordinary norm. The program \path{d_old_compact_secondary.py}
checks all augmentation relations and both residual
generators; their secondary images are zero. The even
endomorphism trace also gives zero divided-orientation
image in both branches. The injection (D.66.9) proves
that both coordinates survive whole compact augmentation.
These computations are the asserted limitations.

\emph{The actual projective edge extension.}
Let \(E=\mathbb Z_2[G]\) be the Bass--Serre edge module and
let \(p:U\to\mathsf M\) be the actual compact projection.
The coefficient-of-identity construction for the projective
target extends \(p\) to an equivariant map
\(a:E\to\mathsf M\). On the selected character fibre,
extend the coefficient functional ordinarily from cycles
to edges, sum its Sylow translates, and average over the
odd remaining cosets. The cycle and edge quotient lattices
are free over \(\mathbb Z_2\), so the ordinary extension
exists. The same formula is integral relative injectivity
and allows the prescribed modulo-four extension. The
scalar average has coefficient \(1/3\), and the
stabilizer relative trace has coefficient
\(1/[H:\operatorname{Syl}_2H]=1/9\).
Modulo four these coefficients are three and one,
exactly as in the program. The
reduction lifting can also be seen from
\(H^1(G,\operatorname{Hom}_{\mathbb Z_2}(E/U,\mathsf M))=0\):
this Hom lattice is a summand of an induced lattice, with
the literal twisting map \(g\otimes v\mapsto g\otimes g^{-1}v\).

Write \(w=a(e_1)\). The actual path cocycle is
\(p_B(s)=0\), \(p_B(r)=e_1-e_r\), whence
\[
 a p_B(s)=0,\qquad a p_B(r)=(1-r)w,\qquad
 a p_B(t)=(s-t)w.
 \tag{D.66.10}
\]
The compact projection kills every actual cusp cycle.
Thus \(N_{36}(t)(s-t)w=0\). The lattice action factors
through the cusp of order twelve in \(G^\dagger\), so
division by three gives the same norm equation there.
Projectivity over the cyclic subgroup gives
\(\ker N_{12}=(t-1)\mathsf M\). Choose
\[
 (t-1)b=(s-t)w,\qquad
 \eta_s=(1-s)b,\qquad
 \eta_r=(1-r)(w+b),\qquad\eta_t=0.
 \tag{D.66.11}
\]
The specified solution modulo four lifts integrally:
the difference of two reduced solutions is \(t\)-fixed;
fixed vectors in a projective cyclic lattice are norm
vectors, whose reductions lift. This proves compatibility
at every integral homology level.

The program \path{d_old_compact_source_cochain.py}
restricts its extension to exactly the retained
96-cycle-to-16-cover map and checks all 72 stabilizer
equations. It induces the full rank-256 real module.
The matrix of \(t-1\) has rank 232. The retained solution
and complete source matrices satisfy (D.66.11) modulo
four exactly.

\emph{A class-two group with the actual exterior map.}
In one real character fibre put
\[
 q(v)=[v\bar v^{\,t}]\in C,\qquad
 c(v,w)=[v\bar w^{\,t}+w\bar v^{\,t}],
 \tag{D.66.12}
\]
where the brackets denote the actual Hermitian coinvariant
map. The polarization of \(q\) is \(c\). Start with the
bilinear map whose diagonal is \(q(e_i)\), whose strict
lower triangle is \(c(e_i,e_j)\), and whose strict upper
triangle is zero. Average its scalar translates by
\(1,\omega,\omega^2\). Since three is one modulo two,
the resulting map \(B\) has the same diagonal and
polarization and is scalar invariant. Sum over the eight
unordered fibres and define
\[
 (v,c)(w,d)=(v+w,c+d+B(\bar v,\bar w)),\qquad
 v,w\in\mathsf M/4,\quad c,d\in C.
 \tag{D.66.13}
\]
Here bars on \(v,w\) mean reduction modulo two.
Bilinearity is the associativity identity. The commutator
is \(c(\bar v,\bar w)\), and the square is
\((2v,q(\bar v))\), so the fourth power is one.
The exterior images span \(C\). Hence the derived group
is \(C\) and the abelianization is \(\mathsf M/4\).

For a deck matrix \(A_g\), the bilinear difference
\(B(A_gv,A_gw)+B(v,w)\) is symmetric with zero diagonal.
Take its triangular quadratic primitive and average
the scalar translates. The resulting function \(k_g\)
gives an automorphism
\[
                     g(v,c)=(A_gv,c+k_g(\bar v)).
 \tag{D.66.14}
\]
The polar identity proves this assertion. Two such
scalar-equivariant corrections differ by an additive
map from a nontrivial odd-character module to a trivial
module, so are equal. This proves all deck composition
identities, including character exchange and permutation
of fibres. The same construction and uniqueness apply
to the actual small map and to any equivariant
automorphism with its induced central map.

In \(\mathscr N\rtimes G^\dagger\) put
\[
 S=\operatorname{Int}((b,0))(s),\qquad
 R=\operatorname{Int}((w+b,0))(r).
 \tag{D.66.15}
\]
Their orders are two and three, as literal conjugates
of the marked finite generators. Their product has
zero abelian coordinate by (D.66.11). The complete
quadratic evaluation gives
\[
                         SR=(0,\mathtt{26};t).
 \tag{D.66.16}
\]
Its eighteenth power has central coordinate zero and
base coordinate \(t^6\). In the actual affine base
this is the cusp involution retained in D.65, with its
separate source marking.

The universal property of the profinite free product
gives a map from \(\widehat B\). On the base kernel
its homology map is the actual onto projection \(p\)
modulo two. Finite Frattini generation makes the image
the whole \(\mathscr N\), proving surjectivity.
Since \(K_3\) maps into a class-two exponent-four
group, the map kills \(L_2\) and factors through \(X\).

The actual \(\phi\) and (D.66.5) similarly give a
unique scalar-equivariant class-two map to the small
group. It commutes with all deck actions by uniqueness,
and is onto by its Frattini map. Together with
\(G^\dagger\to G_*\) it gives \(\mathscr E\to\mathscr E_*\).
The original \(K_*\) maps into \(\mathscr N_*\), so
the composed source map kills \(L_{*,2}\). This proves
the comparison with the actual small class-two source.

\emph{All isometries and the relative quadratic correction.}
The program \path{d_old_compact_all_isometries.py} enumerates
the complete affine spaces of D.64.17 and tests the entire
isometry equations. On this orbit there are six binary
parameters, hence 64 candidates, in each branch. Exactly
16 are isometries in each branch. Relative to the D.64
chosen solution \(X_0\), put \(\delta=X X_0^{-1}\). Then
\[
             \phi\delta=\phi,\qquad
             \delta^t H\bar\delta=H.
 \tag{D.66.17}
\]
The full induced central actions are (D.66.7), each four
times. The functional \(\ell\) is the last coordinate
of (D.66.5); it vanishes on that map's kernel. The
vectors \(\mathtt{09},\mathtt{42}\) are in this kernel.
This proves the claimed pointwise fixedness.

Every such matrix has an actual integral equivariant
lift with \(\phi\delta=\phi\). The reduced
\(\delta-I\) maps to the saturated lattice
\(\ker\phi\), and its domain is projective, so the
equivariant Hom reduction is onto. The certificate
retains a specific lift modulo four: project an
ordinary lift into \(\ker\phi\), then apply the
coefficient-of-identity trace from the projective
domain. Its reduction, all 72 equivariance equations
and its exact small square are checked. The integral
relative average includes \(1/[H:\operatorname{Syl}_2H]=1/9\);
this scalar is one modulo four. Scalar averaging includes
\(1/3\), whose reduction modulo four is three.
Projectivity lifts the specified modulo-four map
integrally. Invertibility modulo two makes it an
integral automorphism. The reduced isometry identity
(D.66.17) is the one used in the present central equation.

Let \(F_\delta\) be its canonical automorphism of
\(\mathscr N\). It commutes with the full deck action.
Its composition with the small class-two map equals
that map: both the linear and central squares commute,
and uniqueness of the scalar-equivariant primitive
proves the full square. Define
\[
 \begin{aligned}
 d_s&=F_\delta((b,0))(b,0)^{-1},\\
 d_r&=F_\delta((w+b,0))(w+b,0)^{-1}.
 \end{aligned}
 \tag{D.66.18}
\]
Both have zero small class-two image, including their
central coordinates. Conjugating (D.66.15) by them
gives its images under \(F_\delta\). Consequently
the literal new cusp differs by
\[
                 (\delta_C-I)c_t=v\ell(c_t)=v.
 \tag{D.66.19}
\]
The programs evaluate both source conjugations and
their product for all 32 cases, with every linear
and quadratic small comparison coefficient. The
source orders remain exact. The eighteenth power
is unchanged because its central difference is
eighteen times an element of order two. This is
the image of the original \(t^{18}\) through the
stated quotient; the full source \(G\) still has
cusp order 36.

\emph{The exact source intersection and integral lift.}
Let \(P=U/\mathcal R\) be actual compactification,
with small cusp lattice \(\mathcal R_*\).
D.64.11 proves the actual onto maps
\[
 \mathcal R\twoheadrightarrow\mathcal R_*,
 \qquad\mathcal L\twoheadrightarrow\ker(P\to P_*).
 \tag{D.66.20}
\]
Their cusp width factor is three, a unit over
\(\mathbb Z_2\). Choose an equivariant affine
section \(\sigma_*\) of the included small compact
summand \(\mathsf M_*\). It exists by projectivity
over its actual small group. Projectivity of
\(\mathsf M\) lifts its compact inclusion to an
affine map. Its difference from \(\sigma_*\phi\)
after the small map lies in \(\mathcal R_*\).
Lift this difference through the first map in
(D.66.20) and subtract it. This gives actual
integral maps
\[
 \begin{gathered}
 \sigma:\mathsf M\longrightarrow U,\qquad
 p\sigma=I,\qquad
 \theta_U\sigma=\sigma_*\phi,\\
                  \sigma(\ker\phi)\subset\mathcal L.
 \end{gathered}
 \tag{D.66.21}
\]
They use the original compact inclusion and
projection. The subtracted cusp term is killed
by \(p\), so these equations hold integrally
with precisely the finite matrices already used.

D.58 identifies the integral homology image of
\(S_{57}\) with \(\mathcal L\). The subgroup
\(R_{50}\cap L_{*,2}\) has image \(\mathcal K\)
in the central layer of \(\mathcal T_2\), by the two actual source
inclusions after D.58.16. D.64's surjective
compact kernel map and (D.66.9) thus give
\[
 \operatorname{im}_{C}(R_{50}\cap L_{*,2})
                           =\ker(C\to C_*).
 \tag{D.66.22}
\]
On these source elements the \(\mathscr N\)-homology
is zero, and the central map is exactly the
calculated exterior map. This proves the
surjectivity with its actual source subgroup.

Equations (D.66.20)--(D.66.22) prove both inclusions
in
\[
        \operatorname{im}_{\mathscr N}(S_{57})
                      =\ker(\mathscr N\to\mathscr N_*).
 \tag{D.66.23}
\]
One follows from \(S_{57}\subset L_{*,2}\) and
the already proved factorization of the small
source map. For the other, lift a prescribed
linear coordinate in \(\ker\phi\) through
(D.66.21) and \(S_{57}\to\mathcal L\).
The remaining discrepancy is central and is
removed by (D.66.22). This also permits a
prescribed full integral homology vector with
the chosen projected linear coordinate.

Use the integral source homology conjugators
\[
 \gamma_s=\sigma((\delta-I)b),\qquad
 \gamma_r=\sigma((\delta-I)(w+b)).
 \tag{D.66.24}
\]
They belong to \(\mathcal L\). Their source
cochain is
\(\eta=\sigma(\delta-I)\eta^{\mathsf M}\),
where \(\eta^{\mathsf M}\) is (D.66.11).
Its value on \(t\) is zero integrally. The
homology map is \(I+\sigma(\delta-I)p\);
it is identity on \(\ker p\) and \(\delta\)
on \(\sigma(\mathsf M)\). Hence it is an
automorphism at every integral level.

Lift (D.66.24) to two elements of \(S_{57}\).
Correct their central coordinates by (D.66.22)
so that their images are exactly (D.66.18).
This retains their entire integral homology.
The resulting source map \(\delta_B\) has
the literal commuting square
\[
             q_E\delta_B=F_\delta q_E:
                              X\longrightarrow\mathscr E.
 \tag{D.66.25}
\]
The equation holds on both marked generators
and thus on the source. No descent of the
preliminary D.65 action to \(\mathscr E\)
is used.

Here is preservation of the exact source
kernel. The two conjugators lie in
\(S_{57}\subset\mathfrak R_h\), so the map
is identity on \(\mathcal W_h\). It is
identity on \(G\), preserves \(K_3\), and
preserves its characteristic homology and
class-two kernels \(M_\infty,L_2\). Through
the \(\mathcal W_h\) map it preserves
\(\mathfrak R_h\). It therefore preserves
\(R_{50}\cap L_2\) and factors through \(X\).
More explicitly, each of the three kernels
\(\mathfrak R_h,M_\infty,L_2\) contains
the defining kernel of \(X\); the intersection
of their three images in \(X\) is trivial.
The source homomorphism therefore kills
their original intersection, without any
assumption of an automorphism of \(\widehat B\).
It is identity on the complete small class-two
source because both conjugators are in
\(L_{*,2}\). Its homology map is invertible,
and its finite class-two map is invertible by
the Frattini criterion. The exterior map
preserves the comparison kernel and is
invertible there. The exact marked
fibre-product description of D.58 now gives
an automorphism of \(X\). It fixes \(t\)
on \(Y\), with a possible error in
\(\mathcal K\) on \(X\).

The normalized parameter comparison is retained
through its actual source map. D.49 gives
\(\mathfrak R_h\subset C_h\subset\widehat F'\).
All correcting conjugators lie in
\(S_{57}\subset\mathfrak R_h\), so their
abelian coordinates in \(\widehat F\) are zero.
The existing common commutator gauge and its
literal image \(f_h^*\) in \(\mathcal W_h\)
are preserved.

\emph{Removal of the two actual residual errors.}
Take a preliminary D.65 lift to \(X\)
with cusp error \(z\in\mathcal J_{65}\),
before quotienting that subgroup. Its
actual conjugators preserve the same source
kernels and give a source automorphism
with that central error. Its image in
\(C\) lies in \(V_0\), by the exact
norm-kernel map following (D.66.9).

The automorphism \(\delta_B\) may have
errors in previously controlled components.
Repair them by the already proved source
procedures of D.59--D.65. These preserve
(D.66.25), as follows. The extra-projective
and cusp corrections have zero
\(\mathsf M\)-homology, hence central
images in \(\mathscr E\); their source
conjugations act trivially there. The
same applies to corrections supported on
the other three character orbits.
Central two-conjugator augmentation
corrections also have zero effect in
\(C\). The changed compact matrix is
an isometry modulo two and is identity
on the orthogonal extra-projective
summands. The actual orientation and
injective-norm arguments therefore restore
the earlier cross and self conditions.
The fourth compact ordinary norm is zero
by the same orientation identity.
These operations put its error in
\(\mathcal J_{65}\) while retaining
(D.66.25). The source-involution words
are followed by the actual moving
augmentation repair, in the order
proved in D.64 and D.65.

Choose \(v=q_E(z)\) in (D.66.7).
The correction preserves \(\mathcal J_{65}\):
the first three compact kernels remain invariant,
it preserves the fourth comparison kernel through
the actual small square \(\phi\delta=\phi\),
and its \(G\)-equivariant exterior map commutes
with the full norm \(\nu_4\).
The repaired correction error and the transported
old error therefore both remain in
\(\mathcal J_{65}\).
The composite error has the exact image
\[
 \begin{aligned}
 q_E\bigl(\delta_B\alpha_{65}(t)t^{-1}\bigr)
    &=\delta_C(q_E(z))+(\delta_C-I)c_t\\
    &=q_E(z)+v=0.
 \end{aligned}
 \tag{D.66.26}
\]
The first equality uses only (D.66.25).
The second uses pointwise fixedness on
the comparison kernel and
\(\ell(c_t)=1\). This is an actual
source error calculation.

The remaining fourth compact error has
zero coinvariant. Apply the actual
augmentation map
\((b,d)\mapsto(1-s)b+s(1-r)d\) to
remove it. The four compact maps from
\(\mathcal J_{63}\) are jointly onto
by D.64, so choose its coefficients
with the other three compact coordinates
zero. Their source lifts are in
\(R_{50}\cap L_{*,2}\). They preserve
full homology, the literal small
conjugators and \(\mathcal W_h\).
The remaining error is exactly in
\(\mathcal J_{66}\).

On \(Y\), the replacement is the
automorphism of (D.66.24), followed
by the D.65 action and the retained
extra-projective cusp-inner repairs:
\[
 \operatorname{Int}(b_{t,Y})\delta_Y\alpha_{65,Y},
 \qquad
 b_t=\prod_j\prod_{i=0}^{35}t^iu_jt^{-i},
 \quad u_j\in S_{57}.
 \tag{D.66.27}
\]
Every one of the 36 factors is retained.
Its \(Y\)-image centralizes \(t\);
its small and \(\mathcal W_h\) images
are zero. The subsequent source-involution
and central augmentation corrections
leave this \(Y\)-action unchanged.
This is the exact comparison with
the previous representative.

The defining kernel \(\mathcal J_{66}\)
is preserved: compact projections kill
cusp terms, the extra left rows are
retained, all four compact reductions
are invertible, and their actual small
squares commute. The finite class-two
map is an automorphism by its Frattini
map. The same literal source conjugators
give all integral homology levels and
retain the prescribed \(\mathcal W_h\)
action. The actual fibre-product maps
therefore prove (D.66.2).

\emph{The remaining kernel and execution scope.}
Joint surjectivity of the compact
kernel maps gives
\(\mathcal J_{65}/\mathcal J_{66}
 \simeq\ker\nu_4\), with the original
source projection. The computed
dimensions give (D.66.3):
\(7682+1918=9600\),
\(4+10+10+4=28\), and
\(276478-1918=274560\).
The positive remaining specified
block proves that (D.66.8) is a
proper nonzero source kernel.

Nine complete programs retain the
actual integral lattices, all norm
and quotient columns, complete
isometry families, the projective
edge extension, full quadratic group
maps, relative small comparisons
and induced norm-kernel module.
Their actual outputs and immutable
input hashes accompany the proof.
The earlier missing-helper and
syntax failures are preserved with
their actual transcripts, as is the
successful norm-kernel run before
its extra central comparison was
added. Only the successful completed
executions support the assertions here.
\end{proof}
